\documentclass[aspectratio=169]{beamer}

%\usepackage[table]{xcolor}
\mode<presentation> {
\setbeamercovered{transparent}
  \usetheme{Boadilla}
\usepackage{booktabs}
\renewcommand{\familydefault}{cmss}

\useinnertheme{rectangles}
}
\usepackage{bm}
\usepackage{amsmath}
\usepackage{stackrel}
\setbeamercolor{normal text}{fg=black}
\setbeamercolor{structure}{fg= blue}
\definecolor{trial}{cmyk}{1,0,0, 0}
\definecolor{trial2}{cmyk}{0.00,0,1, 0}
\definecolor{darkgreen}{rgb}{0,.4, 0.1}
\usepackage{array}
 \usepackage{listings}
\definecolor{darkpurple}{rgb}{0.4, 0, 0.6}
\beamertemplatesolidbackgroundcolor{white}  \setbeamercolor{alerted
text}{fg=darkpurple}
\usepackage{tcolorbox}
\setbeamertemplate{caption}[numbered]\newcounter{mylastframe}

\usepackage{tikz}
\usetikzlibrary{arrows}
\usepackage{colortbl}
\renewcommand{\familydefault}{cmss}

\usecolortheme{lily}

\newtheorem{com}{Comment}
\newtheorem{lem} {Lemma}
\newtheorem{prop}{Proposition}
\newtheorem{condition}{Condition}
\newtheorem{thm}{Theorem}
\newtheorem{defn}{Definition}
\newtheorem{cor}{Corollary}
\newtheorem{obs}{Observation}
 \numberwithin{equation}{section}

\makeatletter
\def\beamerorig@set@color{%
  \pdfliteral{\current@color}%
  \aftergroup\reset@color
}
\def\beamerorig@reset@color{\pdfliteral{\current@color}}
\makeatother

% Define the color and style for the code block
\usepackage{color}
\definecolor{codegreen}{rgb}{0,0.6,0}
\definecolor{codegray}{rgb}{0.5,0.5,0.5}
\definecolor{codepurple}{rgb}{0.58,0,0.82}
\definecolor{backcolour}{rgb}{0.95,0.95,0.92}

\lstdefinestyle{mystyle}{
    backgroundcolor=\color{backcolour},
    commentstyle=\color{codegreen},
    keywordstyle=\color{magenta},
    numberstyle=\tiny\color{codegray},
    stringstyle=\color{codepurple},
    basicstyle=\ttfamily\footnotesize,
    breakatwhitespace=false,
    breaklines=true,
    captionpos=b,
    keepspaces=true,
    numbers=left,
    numbersep=5pt,
    showspaces=false,
    showstringspaces=false,
    showtabs=false,
    tabsize=2
}

\lstset{style=mystyle}
\setbeamertemplate{navigation symbols}{}

\useoutertheme{miniframes}
\title[PLSC 30700]{Week 1 Section: Dummies, Interactions, and Causal Interpretation}

\author{PLSC 30700}
\institute[Chicago]{University of Chicago}
\vspace{0.3in}


\begin{document}


\begin{frame}
\maketitle
\end{frame}


%%%%%%%%%%%%%%%%%%%%%%%%%%%%%%%%%%%%%%%%%%%%%%%%%%%%%%%%%%%%%%%%
\section{Dummy Variables}
%%%%%%%%%%%%%%%%%%%%%%%%%%%%%%%%%%%%%%%%%%%%%%%%%%%%%%%%%%%%%%%%

\setcounter{subsection}{1}
\begin{frame}{Indicator (Dummy) Variables}
\begin{itemize}
\item Our event space $\Omega$ is often binary or categorical (or coded to be so).
\item[-] Examples: Age bins (18--34), Party (R, D), Country (USA, Japan).
\end{itemize}
\bigskip
\begin{defn}
Suppose $A$ is an event.  The \emph{indicator variable} $D = \mathbb{I}(\omega \in A)$ equals 1 if the outcome is in $A$, and 0 otherwise.  Then $\mathbb{E}[D] = P(A)$.
\end{defn}
\medskip
The CEF with a binary regressor is:
$$\mathbb{E}[Y \mid D] = \alpha + \delta D$$
This is the BLP \emph{and} the CEF---the CEF of $Y$ given a binary variable is always linear.
\end{frame}

\begin{frame}{Indicator Variables: Interpretation}

\begin{align*}
\mathbb{E}[Y \mid D = 1] & =\alpha+\delta \\
\mathbb{E}[Y \mid D = 0] & =\alpha \\
\mathbb{E}[Y \mid D = 1]-\mathbb{E}[Y \mid D = 0]&=\delta
\end{align*}
\pause
\begin{itemize}
\item $\alpha$ is the population mean of $Y$ in the reference group ($D=0$).
\item $\delta$ is the population difference in means between the two groups.
\item The BLP slope equals the difference in conditional means---no approximation needed.
\end{itemize}

\end{frame}

\begin{frame}{Parallel Slopes}
\begin{itemize}
\item With a continuous variable $X$ and a group indicator $D$, the \emph{parallel slopes} model is:
$$\mathbb{E}[Y \mid D, X] = \beta_0 + \delta D + \beta_1 X$$
\pause
\item The effect of $X$ is the same for each group: $\partial\, \mathbb{E}[Y|D,X]/\partial X = \beta_1$.
\item The group indicator shifts the intercept by $\delta$ but does not change the slope.
\end{itemize}
\end{frame}


\begin{frame}{Unordered Categorical Variables}
\begin{itemize}
\item Suppose we have multi-category data.
\item For example, ethnic identification $X\in \{ \text{Kazakh, Russian, Other}\}$.  Define:
$$D_{2}=\begin{cases} 1 \text{ if } X=\text{Kazakh} \\ 0 \text{ otherwise}\end{cases} \quad D_{3}=\begin{cases} 1 \text{ if } X=\text{Russian}\\ 0 \text{ otherwise} \end{cases}$$
\item The linear projection becomes:
\begin{align*}
\mathbb{E}[Y \mid D_2, D_3] &= \beta_1 + \beta_2 D_{2} + \beta_3 D_{3}
\end{align*}
\item $\beta_1 = \mathbb{E}[Y \mid \text{Other}]$ (baseline group mean); $\beta_2, \beta_3$ are differences from baseline.

\end{itemize}
\end{frame}

\begin{frame}[fragile]{Dummies and Interactions in R}
\begin{lstlisting}[language=R]
# R handles dummies automatically with factor()
df$party <- factor(df$party)  # R, D, I
mod <- lm(approval ~ party, data = df)
# Intercept = mean for baseline (alphabetically first)
# Coefficients = differences from baseline

# Interaction: does party moderate the economy effect?
mod2 <- lm(approval ~ party * economy, data = df)
# party * economy expands to: party + economy + party:economy

# Marginal effect of economy for Democrats:
coef(mod2)["economy"] + coef(mod2)["partyD:economy"]
\end{lstlisting}
\end{frame}

%%%%%%%%%%%%%%%%%%%%%%%%%%%%%%%%%%%%%%%%%%%%%%%%%%%%%%%%%%%%%%%%
\section{Interaction Terms}
%%%%%%%%%%%%%%%%%%%%%%%%%%%%%%%%%%%%%%%%%%%%%%%%%%%%%%%%%%%%%%%%

\setcounter{subsection}{1}
\begin{frame}{Interactions with Dummy Variables}
\begin{itemize}
\item Suppose now we have a model with different averages and different responses to $x$
\begin{align*}
\text{Group 1: }y_i&=\mu+\beta_1x_i+e_i,\\
\text{Group 2: }y_i&=\mu+\delta+(\beta_1+\gamma)x_i+e_i,\\
\end{align*}
\item Again, we can use a dummy variable $d_i=0$ if $i$ is in group 1, $d_i=1$ if $i$ is in group 2.
\begin{align*}
y_i&=\mu+\delta d_i+ \beta_1x_i+  \gamma d_i x_i+e_i
\end{align*}
\item $d_ix_i$ is called an interaction term.
\end{itemize}
\end{frame}

\begin{frame}{Interactions Generally and Interpretation}
\begin{itemize}
\item With continuous moderator $z$:
\begin{align*}
y_i&=\mu+ \delta z_i +\beta_1 x_i+  \gamma z_i*x_i+e_i
\end{align*} \pause
\item An interaction conditions the effect of $x$ with $z$:
\begin{align*}
\frac{\partial E[y_i|X]}{\partial x_i}&=\beta_1 +  \gamma z_i
\end{align*}
\item A one unit increase in $x_i$ produces a $\beta_1 +  \gamma z_i$ unit increase in $y_i$.
\item $\beta_1$ is the effect when $z_i=0$.
\end{itemize}
\end{frame}

\begin{frame}{The Interpretation Problem with Interactions}
Recall: with an interaction, the marginal effect of $X$ depends on $Z$:
$$\frac{\partial\, \mathbb{E}[Y|X,Z]}{\partial X}=\beta_1 +  \gamma Z$$
\pause
\begin{itemize}
\item $\beta_1$ is the effect of $X$ \emph{when $Z=0$}.
\item But $Z=0$ may not be meaningful!
\item[-] If $Z$ is income, $Z=0$ means zero income---few observations live there.
\item[-] If $Z$ is years of education, $Z=0$ is nonsensical for adults.
\end{itemize}
\bigskip\pause
\textbf{Solution:} Demean $Z$ (and optionally $X$) before interacting.
\end{frame}

\begin{frame}{Demeaning to Get the Average Marginal Effect}
Replace $Z$ with $Z - \mathbb{E}[Z]$ (or $Z - \bar{Z}$ in the sample):
$$\mathbb{E}[Y|X,Z] = \mu + \delta(Z-\mathbb{E}[Z]) + \beta_1 X + \gamma\, (Z-\mathbb{E}[Z])\cdot X$$
\pause
The marginal effect of $X$ is now:
$$\frac{\partial\, \mathbb{E}[Y|X,Z]}{\partial X}=\beta_1 + \gamma(Z - \mathbb{E}[Z])$$
\pause
\begin{itemize}
\item $\beta_1$ is the effect of $X$ \textbf{at the mean of $Z$}---a substantively interpretable quantity.
\item $\gamma$ still measures how the effect of $X$ changes per unit of $Z$.
\item The model is numerically identical---only the interpretation of $\beta_1$ changes.
\end{itemize}
\bigskip\pause
\textbf{Example:} Effect of campaign spending ($X$) moderated by district competitiveness ($Z$).  After demeaning $Z$, $\beta_1$ is the spending effect in a district of average competitiveness.
\end{frame}


%%%%%%%%%%%%%%%%%%%%%%%%%%%%%%%%%%%%%%%%%%%%%%%%%%%%%%%%%%%%%%%%
\section{Causal Interpretation of Regression}
%%%%%%%%%%%%%%%%%%%%%%%%%%%%%%%%%%%%%%%%%%%%%%%%%%%%%%%%%%%%%%%%

\begin{frame}{What Causal Effect Does Regression Recover?}
\begin{itemize}
\item Causal inference is interested in:\\
$$\delta_{X}\equiv E[Y|X, D=1]-E[Y|X, D=0]$$
That is, given some value of X, what is the expected difference in outcomes between treatment and control?
\item Suppose we have the following regression:
$$Y=\delta_R D+X'\beta+e$$
\item It turns out $\delta_R$ is a weighted average of $\delta_X$, where the weights depend on the variance of X.
\item We can use partitions to explain those weights.
\end{itemize}
\end{frame}

\begin{frame}{Applying Partitions to Treatment Dummy}
\begin{align*}
\delta_R&=\frac{Cov(Y, \tilde{D})}{Var (\tilde{D})} \\ \pause
&=\frac{E[(D-E[D|X])Y]}{E[(D-E[D|X])^2]}\\\pause
&=\frac{E\{(D-E[D|X])E[Y|D,X]\}}{E[(D-E[D|X])^2]} \tag{By LIE}\\
\end{align*}
\end{frame}

\begin{frame}{Notational Trick}
\begin{align*}
E[Y|X, D]&=E[Y|X, D=0](1-D)+E[Y|X, D=1]D\\\pause
&=E[Y|X, D=0]+(E[Y|X, D=1]-E[Y|X, D=0])D\\\pause
&=E[Y|X, D=0]+\delta_XD
\end{align*}
\end{frame}

\begin{frame}{Regression as Weighted Average of Causal Effects}
\begin{align*}
&=\frac{E\{(D-E[D|X])E[Y|D,X]\}}{E[(D-E[D|X])^2]}\\\pause
&=\frac{E\{(D-E[D|X])E[Y|X, D=0]\}+E\{(D-E[D|X])D\delta_X\}}{E[(D-E[D|X])^2]} \\\pause
&=\frac{0+E\{(D-E[D|X])D\delta_X\}}{E[(D-E[D|X])^2]} \tag{X and $\tilde{D}$ are uncorrelated.}\\\pause
&=\frac{E[\sigma^2_{D|X}\delta_X]}{E[\sigma^2_{D|X}]} \tag{Where $\sigma^2_{D|X}$ is the conditional variance of D given X.}
\end{align*}
Regression puts the most weight on values of X with the highest variance in D.  Those have the closest to an even split between treatment and control and are most precisely estimated.
\end{frame}


\end{document}
